Eine Kiste 1 (\mEO) liegt auf einem glatten, waagrechten Tisch; die Reibung ist vernachlässigbar. Ein Seil führt von ihr über eine leichte Rolle an der Tischkante zu einer hängenden Kiste 2 (\mZO). Die Kisten werden losgelassen. Rechne mit \gO. Gesucht sind die Beschleunigung und die Seilkraft. Lukas rechnet so:

\begin{center}
\begin{tikzpicture}[>=Stealth]
 \fill[pattern=north east lines] (-0.5,-0.25) rectangle (4,0);
 \draw[thick] (-0.5,0) -- (4,0) -- (4,-2.6);
 \draw[fill=brown!30] (0.4,0) rectangle (1.6,0.8);
 \node at (1,0.4) {1};
 \draw[thick] (4,0) -- (4.3,0.3);
 \draw[thick, fill=white] (4.3,0.3) circle (0.3);
 \fill (4.3,0.3) circle (0.04);
 \draw[thick] (1.6,0.6) -- (4.3,0.6);
 \draw[thick] (4.6,0.3) -- (4.6,-1.2);
 \draw[fill=brown!30] (4.25,-1.2) rectangle (4.95,-1.9);
 \node at (4.6,-1.55) {2};
 % student's free-body diagrams
 \begin{scope}[shift={(7.2,0.2)}]
  \draw[fill=brown!30] (-0.5,0) rectangle (0.5,0.7);
  \node at (0,0.35) {1};
  \draw[->, very thick] (-0.15,0.35) -- (-0.15,-0.9) node[left] {$\vec F_\mathrm{G1}$};
  \draw[->, very thick] (0.15,0) -- (0.15,1.25) node[right] {$\vec F_\mathrm{N1}$};
  \draw[->, very thick] (0.5,0.45) -- (1.35,0.45) node[right] {$\vec F_\mathrm{S}$};
 \end{scope}
 \begin{scope}[shift={(9.4,-0.9)}]
  \draw[fill=brown!30] (-0.35,0) rectangle (0.35,0.7);
  \node at (0,0.35) {2};
  \draw[->, very thick] (0,0.7) -- (0,1.4) node[right] {$\vec F_\mathrm{S}$};
  \draw[->, very thick] (0,0.35) -- (0,-0.6) node[right] {$\vec F_\mathrm{G2}$};
 \end{scope}
\end{tikzpicture}
\end{center}

\begin{enumerate}
\item Kräfte (siehe Zeichnung): auf Kiste 1 Gewichtskraft, Normalkraft und Seilkraft; auf Kiste 2 Gewichtskraft und Seilkraft.
\item Kiste 1, waagrecht (zur Rolle hin positiv): $F_\mathrm{S} = m_1\,a$.
\item Kiste 2: Das Seil hält die Kiste, also ist $F_\mathrm{S} = m_2\,g = \FSP{0}{3}$.
\item Aus 2 und 3: $a = \dfrac{m_2\,g}{m_1} = \aSP{0}{2}$.
\end{enumerate}

\begin{abcliste}
\abc
Welcher Schritt ist falsch? Begründe.

\abc
Stelle die richtige Gleichung auf und berechne die Beschleunigung der Kisten.

\abc
Wie gross ist die Seilkraft wirklich? Vergleiche sie mit der Gewichtskraft von Kiste 2.
\end{abcliste}
